\documentclass[11pt]{article}
\usepackage[T1]{fontenc}
\usepackage[utf8]{inputenc}
\usepackage[brazilian]{babel}
\usepackage{amsmath}
\usepackage{newpxtext,newpxmath}
\usepackage[a4paper,margin=22mm,headheight=15pt,headsep=8mm,footskip=11mm]{geometry}
\usepackage{xcolor,graphicx,booktabs,array,caption,fancyhdr,enumitem}
\usepackage{tikz,pgfplots}
\usepackage[most]{tcolorbox}
\usepackage[hidelinks]{hyperref}
\pgfplotsset{compat=1.18}
\makeatletter\g@addto@macro\UrlBreaks{\do\-}\makeatother
\usetikzlibrary{arrows.meta}

\definecolor{ink}{HTML}{223247}
\definecolor{navy}{HTML}{274568}
\definecolor{quiet}{HTML}{617384}
% Paletas categóricas validadas (dataviz, modo claro): uma por gráfico, para que
% a cor não sugira correspondência entre fonte e geometria.
\definecolor{fOSM}{HTML}{2A78D6}
\definecolor{fAle}{HTML}{EB6834}
\definecolor{fVor}{HTML}{1BAF7A}
\definecolor{gDis}{HTML}{4A3AA7}
\definecolor{gToq}{HTML}{E87BA4}
\definecolor{gSob}{HTML}{008300}

\captionsetup{font=small,labelfont={bf,color=navy},labelsep=colon}
\pagestyle{fancy}\fancyhf{}
\renewcommand{\headrulewidth}{.4pt}
\fancyhead[L]{\small\sffamily\color{ink}SPEEDUP SOBRE O SOLVER DE FEKETE ET AL.}
\fancyhead[R]{\small\sffamily\color{quiet}TouringPolygons}
\fancyfoot[L]{\small\color{quiet}9 de outubro de 2026}
\fancyfoot[R]{\small\color{quiet}\thepage}
\setlength{\parindent}{0pt}\setlength{\parskip}{5pt}
\newcommand{\sect}[1]{\par\medskip{\color{navy}\large\bfseries #1\par}\smallskip}
\newcommand{\vez}{\ensuremath{\times}}

% Opções dos dois eixos (inline: um estilo com ymode é aplicado tarde demais).
\pgfplotsset{pts/.style={only marks,color=#1,mark=*,mark size=1.15pt,mark options={draw=none,fill=#1,fill opacity=.45},forget plot},
  band/.style={#1,line width=1.5pt,mark=*,mark size=2.4pt,mark options={draw=white,line width=.8pt,fill=#1},forget plot}}
\tikzset{endlab/.style={font=\scriptsize\bfseries,text=ink,anchor=west}}

% Dados das figuras, gerados a partir da rodada da dantzig de 08/10
% (float-recovery-dantzig, campanha local não versionada), de
% free-order-dantzig-2026-10-06/per-case.csv e de
% fekete-comparison/analysis/instance-classification.csv. Ver README.md.
% pontos: um caso por linha; x = polígonos com deslocamento
% ((caso*37 mod 11) - 5)*0,11; fonte 0 OSM, 1 aleatórias, 2 Voronoi;
% geom 0 disjuntas, 1 só toques, 2 sobreposição.
% faixas: média geométrica por faixa de 10 polígonos, no número médio de
% polígonos da faixa.
\pgfplotstableread{
caso n x razao fonte geom
1 48 47.890 52.2857 2 1
2 5 5.330 10.2242 0 0
3 5 4.560 21.7123 0 0
4 10 10.000 36.2348 0 0
5 10 10.440 66.6527 0 0
6 15 14.670 75.5362 0 0
7 15 15.110 62.4107 0 0
8 20 20.550 79.9062 0 0
9 20 19.780 102.827 0 0
10 30 30.220 169.732 0 0
11 30 29.450 180.131 0 0
12 40 39.890 169.784 0 0
13 40 40.330 130.489 0 0
14 50 49.560 222.607 0 0
15 50 50.000 298.9 0 0
16 60 60.440 409.071 0 1
17 60 59.670 372.82 0 0
18 5 5.110 23.7134 0 0
19 5 5.550 31.1251 0 0
20 10 9.780 18.1957 0 0
21 10 10.220 38.8946 0 0
22 16 15.450 58.4866 0 1
23 15 14.890 53.5325 0 0
24 20 20.330 137.142 0 2
25 20 19.560 41.445 0 0
26 30 30.000 287.39 0 0
27 30 30.440 229.758 0 0
28 40 39.670 394.591 0 0
29 40 40.110 79.6744 0 1
30 50 50.550 206.087 0 2
31 50 49.780 360.593 0 1
32 60 60.220 277.467 0 1
33 60 59.450 212.977 0 0
34 5 4.890 13.8363 0 0
35 5 5.330 50.6641 0 1
36 10 9.560 16.3678 0 1
37 10 10.000 18.7431 0 1
38 15 15.440 17.888 0 1
39 15 14.670 26.5745 0 1
40 20 20.110 38.1744 0 1
41 20 20.550 108.977 0 1
42 30 29.780 183.839 0 1
43 30 30.220 42.3108 0 1
44 40 39.450 35.9916 0 1
45 40 39.890 104.536 0 1
46 50 50.330 57.6184 0 1
47 50 49.560 167.647 0 1
48 60 60.000 114.202 0 1
49 60 60.440 81.7961 0 2
50 52 51.670 49.9129 2 1
51 5 5.110 17.4447 0 0
52 5 5.550 52.2942 0 1
53 10 9.780 9.83986 0 1
54 10 10.220 24.3395 0 1
55 15 14.450 54.2048 0 1
56 15 14.890 30.8155 0 1
57 20 20.330 144.139 0 1
58 20 19.560 136.385 0 1
59 30 30.000 218.659 0 1
60 30 30.440 179.051 0 0
61 40 39.670 112.685 0 1
62 40 40.110 54.0568 0 1
63 50 50.550 200.023 0 1
64 50 49.780 899.967 0 1
67 5 4.890 14.7511 0 1
68 5 5.330 14.0869 0 1
69 10 9.560 88.9993 0 0
70 6 6.000 24.2845 0 1
71 15 15.440 27.4505 0 1
72 15 14.670 14.4049 0 1
73 20 20.110 66.5952 0 1
74 20 20.550 19.1758 0 1
75 30 29.780 109.127 0 1
76 30 30.220 51.7546 0 1
77 40 39.450 67.5817 0 1
78 40 39.890 1087.35 0 1
79 50 50.330 61.6542 0 1
80 50 49.560 71.7334 0 1
81 60 60.000 147.429 0 1
82 60 60.440 83.201 0 1
83 14 13.670 14.6225 2 1
84 5 5.110 26.1952 0 1
85 5 5.550 30.2143 0 0
86 10 9.780 53.6483 0 0
87 10 10.220 32.3767 0 1
88 15 14.450 59.4976 0 0
89 15 14.890 87.4538 0 1
90 20 20.330 178.134 0 0
91 20 19.560 89.6636 0 1
92 30 30.000 188.194 0 1
93 30 30.440 169.894 0 1
94 40 39.670 329.837 0 1
95 40 40.110 97.7739 0 0
96 50 50.550 587.61 0 0
97 50 49.780 405.804 0 1
98 60 60.220 920.09 0 0
99 60 59.450 926.594 0 1
100 5 4.890 26.4388 0 0
101 5 5.330 70.4149 0 0
102 10 9.560 17.8055 0 0
103 10 10.000 57.703 0 0
104 15 15.440 126.382 0 0
105 15 14.670 30.7942 0 0
106 20 20.110 38.6353 0 0
107 20 20.550 106.332 0 0
108 30 29.780 206.701 0 0
109 30 30.220 139.817 0 0
110 40 39.450 90.7062 0 1
111 40 39.890 162.904 0 0
112 50 50.330 316.898 0 0
113 50 49.560 270.883 0 0
114 60 60.000 386.994 0 1
115 60 60.440 217.013 0 0
116 5 4.670 13.847 0 1
117 5 5.110 14.788 0 0
118 10 10.550 36.9701 0 1
119 9 8.780 59.4909 0 0
120 15 15.220 45.5914 0 1
121 15 14.450 38.1249 0 1
122 20 19.890 107.687 0 1
123 20 20.330 154.246 0 1
124 30 29.560 92.2652 0 1
125 30 30.000 199.71 0 0
126 40 40.440 152.37 0 1
127 40 39.670 192.706 0 1
128 50 50.110 66.1585 0 1
129 50 50.550 195.905 0 1
131 60 60.220 1832.64 0 0
132 51 50.450 37.0356 2 1
133 10 9.890 11.4346 2 1
134 15 15.330 16.4764 2 1
135 20 19.560 27.6938 2 1
136 25 25.000 9.24896 2 1
137 30 30.440 22.1602 2 1
138 35 34.670 21.2306 2 1
139 40 40.110 27.8644 2 1
140 45 45.550 36.3059 2 1
141 50 49.780 35.6496 2 1
142 60 60.220 34.1465 2 1
143 4 3.450 20.9293 0 1
144 5 4.890 24.372 0 0
145 10 10.330 23.5695 0 0
146 10 9.560 24.0997 0 0
147 15 15.000 78.4849 0 1
148 15 15.440 40.5985 0 0
149 20 19.670 35.1325 0 1
150 20 20.110 39.7942 0 0
151 30 30.550 56.5092 0 1
152 30 29.780 94.8231 0 0
153 40 40.220 142.233 0 1
154 40 39.450 184.253 0 0
155 50 49.890 111.487 0 1
156 50 50.330 197.291 0 2
157 59 58.560 275.261 0 2
158 60 60.000 219.652 0 1
159 5 5.440 23.1789 0 0
160 5 4.670 25.529 0 0
161 10 10.110 27.4678 0 0
162 10 10.550 35.1325 0 0
163 15 14.780 86.47 0 0
164 15 15.220 16.7887 0 1
165 20 19.450 146.855 0 1
166 20 19.890 51.5231 0 1
167 30 30.330 76.9802 0 1
168 30 29.560 72.7985 0 1
169 40 40.000 77.0148 0 1
170 40 40.440 135.57 0 0
171 50 49.670 288.489 0 0
172 50 50.110 208.435 0 1
173 60 60.550 262.959 0 1
174 60 59.780 2302.11 0 0
175 5 5.220 35.4359 0 0
176 5 4.450 61.5077 0 1
177 10 9.890 12.6043 0 1
178 10 10.330 14.4372 0 1
179 15 14.560 31.3065 0 1
180 15 15.000 41.7448 0 1
181 20 20.440 23.7964 0 1
182 20 19.670 46.5688 0 1
183 30 30.110 45.8852 0 1
184 30 30.550 128.811 0 0
185 40 39.780 66.5102 0 1
186 40 40.220 162.736 0 0
187 50 49.450 36.959 0 1
188 50 49.890 208.467 0 0
189 60 60.330 180.996 0 1
190 60 59.560 53.9543 0 1
191 4 4.000 63.5919 0 0
192 5 5.440 19.6212 0 0
193 10 9.670 20.7515 0 0
194 10 10.110 12.0584 0 0
195 15 15.550 41.7992 0 1
196 15 14.780 44.8303 0 0
197 20 20.220 115.46 0 0
198 20 19.450 60.1287 0 0
199 30 29.890 185.65 0 0
200 30 30.330 71.6149 0 0
201 40 39.560 145.06 0 0
202 40 40.000 165.469 0 0
203 50 50.440 686.674 0 0
204 50 49.670 138.856 0 0
205 60 60.110 247.849 0 0
206 60 60.550 635.293 0 0
207 10 9.780 18.3545 2 1
208 15 15.220 23.3608 2 1
209 20 19.450 15.2286 2 1
210 25 24.890 20.5931 2 1
211 30 30.330 21.2175 2 1
212 35 34.560 12.3796 2 1
213 40 40.000 22.9849 2 1
214 45 45.440 47.2064 2 1
215 50 49.670 13.3723 2 1
216 60 60.110 105.252 2 1
217 4 4.550 15.6138 0 0
218 5 4.780 13.8197 0 0
219 10 10.220 17.6555 0 0
220 10 9.450 26.237 0 0
221 15 14.890 35.9877 0 0
222 15 15.330 52.9173 0 0
223 20 19.560 36.2842 0 0
224 20 20.000 118.358 0 0
225 30 30.440 34.5381 0 1
226 30 29.670 87.7148 0 1
227 40 40.110 82.4488 0 1
228 40 40.550 136.247 0 1
229 50 49.780 182.052 0 1
230 50 50.220 213.69 0 1
231 60 59.450 2010.89 0 1
232 60 59.890 102.417 0 1
233 5 5.330 21.4352 0 1
234 4 3.560 14.1909 0 0
235 10 10.000 20.532 0 0
236 10 10.440 39.7224 0 0
237 15 14.670 63.2057 0 0
238 15 15.110 60.4634 0 0
239 20 20.550 11.4164 0 1
240 20 19.780 40.9067 0 0
241 30 30.220 142.045 0 0
242 30 29.450 100.03 0 0
243 40 39.890 161.374 0 0
244 40 40.330 122.175 0 0
245 50 49.560 404.792 0 0
246 50 50.000 323.124 0 0
247 60 60.440 749.497 0 0
248 60 59.670 1232.44 0 0
249 5 5.110 23.3326 0 0
250 5 5.550 30.2465 0 0
251 10 9.780 38.0936 0 0
252 10 10.220 46.4477 0 1
253 15 14.450 64.159 0 0
254 15 14.890 174.313 0 0
255 20 20.330 78.4357 0 0
256 20 19.560 154.563 0 0
257 30 30.000 231.575 0 0
258 30 30.440 203.953 0 0
259 40 39.670 204.038 0 1
260 40 40.110 177.335 0 1
261 50 50.550 471.729 0 0
262 50 49.780 339.646 0 1
263 60 60.220 540.252 0 1
264 60 59.450 447.479 0 0
265 5 4.890 33.3434 1 0
266 5 5.330 30.4652 1 0
267 5 4.560 38.2785 1 0
268 5 5.000 40.2318 1 0
269 5 5.440 42.5373 1 0
270 5 4.670 68.7018 1 0
271 5 5.110 43.7661 1 0
272 5 5.550 44.1868 1 0
273 5 4.780 20.34 1 0
274 5 5.220 63.4614 1 0
275 5 4.450 30.4168 1 0
276 5 4.890 35.8182 1 0
277 5 5.330 26.5172 1 0
278 5 4.560 56.4337 1 0
279 5 5.000 57.9389 1 0
280 5 5.440 40.5932 1 0
281 5 4.670 29.9215 1 2
282 5 5.110 26.7858 1 2
283 5 5.550 41.8514 1 0
284 5 4.780 56.4299 1 0
285 10 10.220 104.85 1 2
286 10 9.450 81.5328 1 2
287 10 9.890 42.9702 1 0
288 9 9.330 43.6915 1 0
289 10 9.560 29.7956 1 2
290 10 10.000 24.5603 1 2
291 10 10.440 38.1168 1 2
292 10 9.670 41.994 1 0
293 10 10.110 39.8487 1 2
294 10 10.550 29.5707 1 2
295 10 9.780 28.7702 1 2
296 10 10.220 44.8351 1 2
297 10 9.450 39.7775 1 2
298 10 9.890 33.348 1 2
299 10 10.330 27.2369 1 2
300 10 9.560 41.7131 1 2
301 10 10.000 36.7055 1 2
302 10 10.440 32.14 1 2
303 10 9.670 34.7582 1 2
304 10 10.110 41.4391 1 2
305 15 15.550 36.7012 1 2
306 15 14.780 45.8929 1 2
307 15 15.220 50.3027 1 2
308 15 14.450 38.257 1 2
309 14 13.890 84.2002 1 2
310 15 15.330 51.3303 1 2
311 15 14.560 60.6402 1 2
312 15 15.000 111.399 1 2
313 15 15.440 27.8532 1 2
314 15 14.670 46.3049 1 2
315 15 15.110 61.4481 1 2
316 15 15.550 60.0452 1 2
317 15 14.780 55.5096 1 2
318 15 15.220 37.0028 1 2
319 15 14.450 31.5519 1 2
320 15 14.890 63.2548 1 2
321 15 15.330 36.8244 1 2
322 15 14.560 54.2235 1 2
323 15 15.000 61.5808 1 2
324 15 15.440 34.3316 1 2
325 20 19.670 54.0552 1 2
326 20 20.110 92.627 1 2
327 20 20.550 49.519 1 2
328 20 19.780 39.7895 1 2
329 20 20.220 94.1484 1 2
330 20 19.450 34.0284 1 2
331 20 19.890 47.3862 1 2
332 19 19.330 105.89 1 2
333 19 18.560 43.4539 1 2
334 19 19.000 87.5531 1 2
335 19 19.440 115.798 1 2
336 20 19.670 93.6956 1 2
337 20 20.110 100.744 1 2
338 19 19.550 63.4631 1 2
339 20 19.780 63.3679 1 2
340 20 20.220 107.197 1 2
341 19 18.450 187.542 1 2
342 20 19.890 81.8223 1 2
343 20 20.330 73.558 1 2
344 20 19.560 36.7234 1 2
345 30 30.000 84.2193 1 2
346 30 30.440 39.9779 1 2
347 29 28.670 63.9536 1 2
348 30 30.110 71.5428 1 2
349 30 30.550 80.0158 1 2
350 30 29.780 92.138 1 2
351 30 30.220 83.9787 1 2
352 29 28.450 70.7154 1 2
353 30 29.890 163.712 1 2
354 28 28.330 39.6879 1 2
355 30 29.560 124.442 1 2
356 30 30.000 54.3789 1 2
357 30 30.440 85.2389 1 2
358 30 29.670 70.8229 1 2
359 30 30.110 106.377 1 2
360 30 30.550 85.3585 1 2
361 30 29.780 77.9687 1 2
362 30 30.220 37.9813 1 2
363 29 28.450 96.4904 1 2
364 30 29.890 34.8634 1 2
365 40 40.330 68.2548 1 2
366 40 39.560 60.2232 1 2
367 39 39.000 42.8554 1 2
368 40 40.440 36.3263 1 2
369 38 37.670 113.209 1 2
370 38 38.110 63.2137 1 2
371 39 39.550 46.0942 1 2
372 39 38.780 102.21 1 2
373 39 39.220 110.495 1 2
374 38 37.450 79.8646 1 2
375 39 38.890 48.4283 1 2
376 37 37.330 63.3472 1 2
377 39 38.560 61.1767 1 2
378 40 40.000 52.5412 1 2
379 39 39.440 29.5986 1 2
380 40 39.670 51.7571 1 2
381 40 40.110 44.3121 1 2
382 40 40.550 63.6333 1 2
383 39 38.780 50.9385 1 2
384 39 39.220 111.467 1 2
385 48 47.450 130.217 1 2
386 50 49.890 47.7639 1 2
387 48 48.330 138.959 1 2
388 47 46.560 134.35 1 2
389 50 50.000 92.862 1 2
390 49 49.440 173.693 1 2
391 50 49.670 94.2345 1 2
392 48 48.110 239.404 1 2
393 46 46.550 101.641 1 2
394 49 48.780 115.261 1 2
395 49 49.220 109.784 1 2
396 49 48.450 104.432 1 2
397 49 48.890 150.873 1 2
398 48 48.330 78.8947 1 2
399 49 48.560 108.901 1 2
400 50 50.000 48.4609 1 2
401 50 50.440 80.263 1 2
402 48 47.670 208.25 1 2
403 50 50.110 100.432 1 2
404 49 49.550 60.2273 1 2
405 60 59.780 121.771 1 2
406 59 59.220 114.414 1 2
407 60 59.450 81.3643 1 2
408 58 57.890 86.5602 1 2
409 56 56.330 56.0815 1 2
410 58 57.560 140.988 1 2
411 59 59.000 206.79 1 2
412 57 57.440 130.922 1 2
413 59 58.670 86.5994 1 2
414 58 58.110 95.7542 1 2
415 59 59.550 149.271 1 2
416 60 59.780 127.943 1 2
417 59 59.220 143.116 1 2
418 59 58.450 860.923 1 2
419 59 58.890 146.279 1 2
420 60 60.330 549.048 1 2
421 58 57.560 186.675 1 2
422 58 58.000 532.989 1 2
423 59 59.440 106.693 1 2
424 56 55.670 184.811 1 2
425 4 4.110 14.5966 0 1
426 4 4.550 42.9221 0 1
427 10 9.780 9.75668 0 1
428 10 10.220 16.1711 0 1
429 15 14.450 12.9318 0 1
430 15 14.890 13.1672 0 1
431 20 20.330 18.2035 0 1
432 20 19.560 27.4166 0 1
433 30 30.000 70.3299 0 1
434 30 30.440 47.3635 0 2
435 40 39.670 53.6517 0 1
436 39 39.110 81.9963 0 1
437 50 50.550 29.3798 0 2
438 50 49.780 99.0764 0 1
439 59 59.220 87.8948 0 1
440 59 58.450 41.9513 0 1
441 10 9.890 11.3824 2 1
442 20 20.330 9.003 2 1
443 30 29.560 32.8636 2 1
444 40 40.000 35.4966 2 1
445 50 50.440 21.9568 2 1
446 60 59.670 25.1986 2 1
447 10 10.110 19.8277 2 1
448 20 20.550 15.5422 2 1
449 30 29.780 32.9666 2 1
450 40 40.220 25.9135 2 1
451 50 49.450 31.298 2 1
452 60 59.890 8.69921 2 1
453 10 10.330 14.9709 2 1
454 15 14.560 12.6434 2 1
455 20 20.000 23.7729 2 1
456 25 25.440 19.5135 2 1
457 30 29.670 22.7615 2 1
458 35 35.110 22.0903 2 1
459 40 40.550 28.9934 2 1
460 45 44.780 28.9329 2 1
461 50 50.220 26.314 2 1
462 60 59.450 47.1017 2 1
463 4 3.890 38.6462 0 0
464 5 5.330 43.8778 0 0
465 10 9.560 22.1666 0 0
466 10 10.000 25.9685 0 0
467 15 15.440 49.3394 0 0
468 15 14.670 39.2821 0 0
469 20 20.110 84.2995 0 0
470 20 20.550 86.0328 0 0
471 30 29.780 246.949 0 0
472 30 30.220 80.7529 0 2
473 40 39.450 128.209 0 0
474 40 39.890 294.743 0 2
475 50 50.330 380.55 0 0
476 50 49.560 505.406 0 1
477 60 60.000 1665.26 0 0
478 60 60.440 415.956 0 0
479 5 4.670 99.9273 0 0
480 5 5.110 42.7584 0 0
481 10 10.550 12.5442 0 0
482 10 9.780 22.404 0 0
483 15 15.220 59.697 0 0
484 15 14.450 77.6056 0 0
485 20 19.890 51.6645 0 0
486 20 20.330 19.499 0 0
487 30 29.560 138.716 0 0
488 30 30.000 255.438 0 0
489 40 40.440 212.877 0 0
490 40 39.670 151.382 0 0
491 50 50.110 234.562 0 0
492 50 50.550 719.359 0 0
494 60 60.220 728.585 0 1
495 16 15.450 33.6626 2 1
496 22 21.890 41.602 2 1
497 10 10.330 32.5553 2 1
498 10 9.560 28.885 2 1
499 15 15.000 15.8571 2 1
500 15 15.440 36.224 2 1
501 20 19.670 27.7478 2 1
502 20 20.110 24.4239 2 1
503 25 25.550 33.7966 2 1
504 25 24.780 25.625 2 1
505 30 30.220 30.432 2 1
506 30 29.450 38.951 2 1
507 35 34.890 21.4404 2 1
508 35 35.330 11.6639 2 1
509 40 39.560 26.2702 2 1
510 40 40.000 20.47 2 1
511 45 45.440 19.9931 2 1
512 45 44.670 23.5812 2 1
513 50 50.110 45.4677 2 1
514 50 50.550 40.1027 2 1
515 60 59.780 37.6076 2 1
516 60 60.220 80.8547 2 1
517 10 9.450 27.1813 2 1
518 15 14.890 21.908 2 1
519 20 20.330 43.2695 2 1
520 25 24.560 32.5267 2 1
521 30 30.000 24.5033 2 1
522 35 35.440 44.387 2 1
523 40 39.670 42.2051 2 1
524 45 45.110 23.5704 2 1
525 50 50.550 26.1034 2 1
526 60 59.780 31.5835 2 1
527 5 5.220 16.8439 0 0
528 4 3.450 25.0895 0 0
529 10 9.890 56.4711 0 1
530 10 10.330 50.8296 0 0
531 15 14.560 87.8903 0 0
532 15 15.000 38.7587 0 1
533 20 20.440 80.9915 0 1
534 20 19.670 36.4375 0 1
535 30 30.110 182.254 0 0
536 30 30.550 122.333 0 1
537 40 39.780 40.6229 0 1
538 40 40.220 91.1195 0 1
539 50 49.450 259.12 0 1
540 50 49.890 252.82 0 1
541 60 60.330 115.793 0 1
543 5 5.000 22.5597 0 1
544 5 5.440 29.0115 0 1
545 10 9.670 14.953 0 0
546 10 10.110 42.5184 0 0
547 15 15.550 88.9482 0 0
548 15 14.780 52.8826 0 0
549 20 20.220 82.3305 0 1
550 20 19.450 209.392 0 1
551 30 29.890 212.979 0 1
552 30 30.330 147.03 0 1
553 40 39.560 142.893 0 1
554 40 40.000 48.1234 0 1
555 50 50.440 158.782 0 1
556 50 49.670 250.696 0 1
557 60 60.110 696.615 0 1
558 60 60.550 1945.86 0 1
}\dadospontos
\pgfplotstableread{
n g casos
7.34 26.1626 80
17.51 54.4138 80
30.00 123.5127 40
39.98 126.6596 40
50.00 214.9581 40
59.91 349.2660 35
}\dadosfontea
\pgfplotstableread{
n g casos
6.91 30.9172 75
17.25 64.8402 40
30.00 169.0327 20
40.00 159.5644 15
50.00 335.3384 15
60.00 639.8880 13
}\dadosgeoma
\pgfplotstableread{
n g casos
7.47 39.3155 40
17.32 59.4397 40
29.75 72.5216 20
39.10 60.6960 20
48.80 106.6206 20
58.55 155.0618 20
}\dadosfonteb
\pgfplotstableread{
n g casos
7.88 22.2642 34
17.65 35.6142 55
28.85 52.2353 33
39.18 61.5464 38
49.14 85.0999 37
59.37 135.9083 30
}\dadosgeomb
\pgfplotstableread{
n g casos
10.00 19.1551 8
17.50 20.8133 16
27.47 25.7686 15
37.86 24.2967 14
47.87 29.6150 15
58.30 38.1834 10
}\dadosfontec
\pgfplotstableread{
n g casos
9.47 37.3409 19
17.39 60.6642 41
29.77 71.4792 22
39.14 65.4396 21
48.96 106.5532 23
58.64 154.5994 22
}\dadosgeomc


\begin{document}

{\color{navy}\LARGE\bfseries Speedup sobre o solver de Fekete et al.\par}
\vspace{2pt}
{\color{quiet}\large por fonte e por geometria das instâncias, e onde o polimento é usado\par}
\vspace{10pt}

\begin{tcolorbox}[colback=navy!4,colframe=navy!35,boxrule=.5pt,arc=2pt,left=8pt,right=8pt,top=5pt,bottom=5pt]
\small Nos 553 casos do corpus de Fekete et al.\ que o solver deles fecha, o nosso foi mais rápido em todos, com razão mínima de 8,7\vez{} e média geométrica de 65,4\vez. A vantagem cresce com o número de polígonos. Por fonte, a média geométrica é de 81,7\vez{} nas instâncias OSM, 66,8\vez{} nas aleatórias e 25,6\vez{} nas Voronoi. Por geometria, 77,4\vez{} nas disjuntas, 74,1\vez{} nas com sobreposição de área e 52,9\vez{} nas que só se tocam; esta última média é puxada para baixo pelas Voronoi, já que as OSM que só se tocam ficam em 77,4\vez. Das 27,5 milhões de chamadas ao oráculo convexo, 1,38\% foram ao polimento por pontos interiores, cerca de um terço de cada fonte, quase todas em instâncias com toques ou sobreposição.
\end{tcolorbox}

\sect{1\quad Speedup por fonte das instâncias}
\begin{figure}[ht]\centering
\begin{tikzpicture}
\begin{axis}[width=.97\linewidth,height=88mm,ymode=log,log basis y=10,
  xmin=3,xmax=66,ymin=.8,ymax=4000,
  ytick={1,2,5,10,20,50,100,200,500,1000,2000},
  yticklabels={1\vez,2\vez,5\vez,10\vez,20\vez,50\vez,100\vez,200\vez,500\vez,1.000\vez,2.000\vez},
  xtick={5,10,...,60},xlabel={número de polígonos},ylabel={razão de tempo Fekete / nosso},
  label style={font=\small},tick label style={font=\scriptsize},
  axis lines*=left,ymajorgrids,grid style={quiet!18},
  legend style={font=\scriptsize,draw=none,at={(.5,-.17)},anchor=north,legend columns=4,
    /tikz/every even column/.append style={column sep=6pt}},
  every axis plot/.append style={line join=round}]
\addlegendimage{only marks,color=fOSM,mark=*,mark size=2pt,mark options={draw=none,fill=fOSM}}\addlegendentry{OSM}
\addlegendimage{only marks,color=fAle,mark=*,mark size=2pt,mark options={draw=none,fill=fAle}}\addlegendentry{Aleatórias}
\addlegendimage{only marks,color=fVor,mark=*,mark size=2pt,mark options={draw=none,fill=fVor}}\addlegendentry{Voronoi / tessellation}
\addlegendimage{ink,line width=1.5pt}\addlegendentry{média geométrica por faixa}
\addplot[pts=fOSM] table[x=x,y=razao,restrict expr to domain={\thisrow{fonte}}{0:0}] {\dadospontos};
\addplot[pts=fAle] table[x=x,y=razao,restrict expr to domain={\thisrow{fonte}}{1:1}] {\dadospontos};
\addplot[pts=fVor] table[x=x,y=razao,restrict expr to domain={\thisrow{fonte}}{2:2}] {\dadospontos};
\addplot[band=fOSM] table[x=n,y=g] {\dadosfontea};
\addplot[band=fAle] table[x=n,y=g] {\dadosfonteb};
\addplot[band=fVor] table[x=n,y=g] {\dadosfontec};
\addplot[quiet,dashed,line width=.6pt,forget plot] coordinates {(3,1) (66,1)};
\node[font=\scriptsize,text=quiet,anchor=south east] at (axis cs:66,1) {empate, 1\vez};
\node[endlab] at (axis cs:61.4,349.3) {349\vez};
\node[endlab] at (axis cs:61.4,155.1) {155\vez};
\node[endlab] at (axis cs:61.4,38.2) {38,2\vez};
\end{axis}
\end{tikzpicture}
\caption{Razão entre o tempo do solver de Fekete et al.\ e o nosso, por caso (pontos, com pequeno deslocamento horizontal para não se sobreporem) e média geométrica por faixa de 10 polígonos (linhas, no número médio de polígonos da faixa). Escala vertical logarítmica; acima de 1\vez{} o nosso é mais rápido. 553 casos.}\label{fig:fonte}
\end{figure}

\begin{table}[ht]\centering\footnotesize\setlength{\tabcolsep}{4.2pt}
\begin{tabular}{@{}lrrrrrrrr@{}}\toprule
& & \multicolumn{4}{c}{Razão Fekete / nosso} & \multicolumn{3}{c}{Instâncias por geometria}\\\cmidrule(lr){3-6}\cmidrule(l){7-9}
Fonte & Casos & média geom. & mediana & mínimo & máximo & disjuntas & só toques & sobreposição\\\midrule
OSM & 315 & 81,7\vez & 78,5\vez & 9,8\vez & 2.302\vez & 158 & 153 & 9\\
Aleatórias & 160 & 66,8\vez & 62,4\vez & 20,3\vez & 861\vez & 21 & 0 & 139\\
Voronoi / tessellation & 78 & 25,6\vez & 26,0\vez & 8,7\vez & 105\vez & 0 & 78 & 0\\\midrule
Total & 553 & 65,4\vez & 57,6\vez & 8,7\vez & 2.302\vez & 179 & 231 & 148\\\bottomrule
\end{tabular}
\caption{Speedup por fonte nos casos que os dois solvers fecham. As colunas de geometria contam todas as 558 instâncias, incluindo as 5 OSM que o solver de Fekete et al.\ não fechou.}\label{tab:fonte}
\end{table}

\clearpage
\sect{2\quad Speedup por geometria}
A geometria vem dos polígonos originais de cada instância, classificados com Shapely: \emph{disjuntas} não têm nenhum par de polígonos que se toque; \emph{só toques} têm pares que compartilham fronteira, sem área em comum; \emph{sobreposição} têm ao menos um par com área em comum.

\begin{figure}[ht]\centering
\begin{tikzpicture}
\begin{axis}[width=.97\linewidth,height=88mm,ymode=log,log basis y=10,
  xmin=3,xmax=66,ymin=.8,ymax=4000,
  ytick={1,2,5,10,20,50,100,200,500,1000,2000},
  yticklabels={1\vez,2\vez,5\vez,10\vez,20\vez,50\vez,100\vez,200\vez,500\vez,1.000\vez,2.000\vez},
  xtick={5,10,...,60},xlabel={número de polígonos},ylabel={razão de tempo Fekete / nosso},
  label style={font=\small},tick label style={font=\scriptsize},
  axis lines*=left,ymajorgrids,grid style={quiet!18},
  legend style={font=\scriptsize,draw=none,at={(.5,-.17)},anchor=north,legend columns=4,
    /tikz/every even column/.append style={column sep=6pt}},
  every axis plot/.append style={line join=round}]
\addlegendimage{only marks,color=gDis,mark=*,mark size=2pt,mark options={draw=none,fill=gDis}}\addlegendentry{Disjuntas}
\addlegendimage{only marks,color=gToq,mark=*,mark size=2pt,mark options={draw=none,fill=gToq}}\addlegendentry{Só toques de fronteira}
\addlegendimage{only marks,color=gSob,mark=*,mark size=2pt,mark options={draw=none,fill=gSob}}\addlegendentry{Sobreposição com área}
\addlegendimage{ink,line width=1.5pt}\addlegendentry{média geométrica por faixa}
\addplot[pts=gDis] table[x=x,y=razao,restrict expr to domain={\thisrow{geom}}{0:0}] {\dadospontos};
\addplot[pts=gToq] table[x=x,y=razao,restrict expr to domain={\thisrow{geom}}{1:1}] {\dadospontos};
\addplot[pts=gSob] table[x=x,y=razao,restrict expr to domain={\thisrow{geom}}{2:2}] {\dadospontos};
\addplot[band=gDis] table[x=n,y=g] {\dadosgeoma};
\addplot[band=gToq] table[x=n,y=g] {\dadosgeomb};
\addplot[band=gSob] table[x=n,y=g] {\dadosgeomc};
\addplot[quiet,dashed,line width=.6pt,forget plot] coordinates {(3,1) (66,1)};
\node[font=\scriptsize,text=quiet,anchor=south east] at (axis cs:66,1) {empate, 1\vez};
\node[endlab] at (axis cs:61.4,639.9) {640\vez};
\node[endlab,yshift=5pt] at (axis cs:61.4,154.6) {155\vez};
\node[endlab,yshift=-5pt] at (axis cs:61.4,135.9) {136\vez};
\end{axis}
\end{tikzpicture}
\caption{O mesmo gráfico da Figura~\ref{fig:fonte}, com as instâncias agrupadas pela geometria dos polígonos originais. 553 casos.}\label{fig:geometria}
\end{figure}

\begin{table}[ht]\centering\footnotesize\setlength{\tabcolsep}{4.2pt}
\begin{tabular}{@{}lrrrrrrrr@{}}\toprule
& & \multicolumn{4}{c}{Razão Fekete / nosso} & \multicolumn{3}{c}{Instâncias por fonte}\\\cmidrule(lr){3-6}\cmidrule(l){7-9}
Geometria & Casos & média geom. & mediana & mínimo & máximo & OSM & aleatórias & Voronoi\\\midrule
Disjuntas & 178 & 77,4\vez & 62,8\vez & 10,2\vez & 2.302\vez & 158 & 21 & 0\\
Só toques de fronteira & 227 & 52,9\vez & 42,0\vez & 8,7\vez & 2.011\vez & 153 & 0 & 78\\
Sobreposição com área & 148 & 74,1\vez & 72,6\vez & 24,6\vez & 861\vez & 9 & 139 & 0\\\bottomrule
\end{tabular}
\caption{Speedup por geometria nos casos que os dois solvers fecham. As colunas de fonte contam todas as 558 instâncias.}\label{tab:geometria}
\end{table}

A geometria sozinha não explica a diferença entre as classes. Dentro das instâncias que só se tocam, as 149 OSM fechadas pelos dois têm média geométrica de 77,4\vez, perto das OSM disjuntas (84,2\vez, 157 casos); o que puxa a média do grupo para baixo são as 78 Voronoi (25,6\vez). As aleatórias disjuntas, só 21 casos e todos pequenos (5,7 polígonos em média), ficam em 41,2\vez.

\clearpage
\sect{3\quad Onde o polimento por pontos interiores é usado}
Na rodada completa, o oráculo convexo foi chamado 27.525.200 vezes nas 558 instâncias. A prova intervalar fechou 27.145.795 chamadas, e as outras 379.405 (1,38\%) foram ao polimento por pontos interiores, que fechou todas; nenhuma chamada precisou do fluxo exato, em aritmética racional.

\begin{table}[ht]\centering\footnotesize
\begin{tabular}{@{}lrrrrr@{}}\toprule
& Casos & Chamadas & Polimento & \% das chamadas da classe & Parcela do polimento\\\midrule
\multicolumn{6}{@{}l}{\itshape Por fonte}\\
OSM & 320 & 26.199.008 & 119.919 & 0,46\% & 31,6\%\\
Aleatórias & 160 & 1.042.580 & 136.470 & 13,09\% & 36,0\%\\
Voronoi / tessellation & 78 & 283.612 & 123.016 & 43,37\% & 32,4\%\\\midrule
\multicolumn{6}{@{}l}{\itshape Por geometria}\\
Disjuntas & 179 & 11.560.574 & 7 & 0,00006\% & 0,002\%\\
Só toques de fronteira & 231 & 14.841.127 & 239.704 & 1,62\% & 63,2\%\\
Sobreposição com área & 148 & 1.123.499 & 139.694 & 12,43\% & 36,8\%\\\midrule
Total & 558 & 27.525.200 & 379.405 & 1,38\% & 100\%\\\bottomrule
\end{tabular}
\caption{Chamadas ao oráculo convexo e ao polimento, por classe. ``\% das chamadas da classe'': fração das chamadas daquela classe que foram ao polimento; ``parcela do polimento'': fração do total de 379.405.}\label{tab:polimento}
\end{table}

\begin{itemize}[leftmargin=*,itemsep=2pt]
\item \textbf{Por fonte, cerca de um terço cada.} O peso dentro de cada fonte, porém, é muito diferente: as OSM fazem 95\% de todas as chamadas, e quase todas fecham na prova intervalar; nas Voronoi, 43\% das chamadas vão ao polimento.
\item \textbf{A geometria decide.} Nas 179 instâncias disjuntas, só 7 chamadas foram ao polimento. Com polígonos disjuntos, a construção em \texttt{double} quase sempre acerta e a prova intervalar fecha; é nos toques e nas sobreposições que aparecem os empates exatos e os contatos coincidentes.
\item \textbf{As chamadas se concentram.} Dez casos somam 75\% das chamadas ao polimento; só o caso 420 (aleatório, com sobreposição) tem 22\%.
\end{itemize}

\sect{4\quad Dados e limitações}
\begin{itemize}[leftmargin=*,itemsep=2pt]\small\sloppy
\item \textbf{Problema.} TPP de ordem livre com extremos fixos, nas 558 instâncias do corpus de Fekete et al.\ (\nolinkurl{benchmarks/results-saved/fekete-comparison/instances.bin}). Os dois solvers param com gap relativo de 0,1\% (\(U\le1{,}001\,L\)), uma thread por caso. ``Fechou'' quer dizer gap fechado, não ótimo algébrico.
\item \textbf{Nosso solver.} Padrão atual (prova intervalar e polimento em ponto flutuante, com o fluxo exato como reserva), rodada de 08/10 na dantzig (Intel i9-12900K), binário \texttt{ded35f2}, 4 casos simultâneos com a máquina compartilhada (\emph{load average} em torno de 7). Fechou e validou os 558 casos (Shapely, tolerância \(10^{-7}\)). As contagens de chamadas são dessa mesma rodada.
\item \textbf{Solver de Fekete et al.} Configuração original (\texttt{FEASIBILITY\_TOLERANCE} \(10^{-3}\)), campanha de 02/10 na mesma máquina, com vários casos em paralelo; três casos (131, 420 e 558) vêm de uma execução posterior. Fechou 553 casos; os 5 restantes (65, 66, 130, 493 e 542, todos OSM) ficaram dias sem fechar e estão fora dos gráficos e das razões, o que subestima a vantagem. Os caminhos dessa campanha não têm validação independente registrada.
\item \textbf{Comparação.} Os tempos não vêm de uma rodada lado a lado; os casos de poucos milissegundos são os mais sensíveis a carga e a diferenças de ambiente. Uma afirmação para publicação exige rodar os dois solvers juntos, no mesmo ambiente.
\item \textbf{Fontes dos dados} (em \nolinkurl{benchmarks/results-saved}, salvo indicação). Tempos do Fekete: \nolinkurl{free-order-dantzig-2026-10-06/per-case.csv}. Fonte e geometria: \nolinkurl{fekete-comparison/analysis/instance-classification.csv}. Nossos tempos e chamadas: campanha local \nolinkurl{float-recovery-dantzig}, resumida na seção \nolinkurl{float-recovery-dantzig-2026-10-08} do \nolinkurl{README.md}.
\end{itemize}

\end{document}
